\documentclass[lualatex]{article}

\usepackage{ctex}
\usepackage{amsmath}
\usepackage{amssymb}
\usepackage{mathtools}
\usepackage{bm}
\usepackage[margin=1in]{geometry}

\title{
    从 RCD 空间到 LCD 空间\\[0.5em]
    \large 低正则性理论迷失的二十年和今后的二十年
}

\author{Arnold Robinson}
\date{\today}

\begin{document}

\maketitle

\section{Introduction}

一切源自于 Poincar\'e 在 1904 年提出的著名问题：

\begin{quote}
所有单连通三维闭流形是否与三维球面同胚？
\end{quote}

这个问题极大地改变了 20 世纪几何学的发展，并激发了无数强力的工具和理论。Hamilton 在 1980 年代提出的 Ricci flow
\[
    \frac{\partial}{\partial t} g_t = -2\operatorname{Ric}(g_t)
\]
提供了一种完全不同于传统拓扑方法的视角：让几何结构本身随时间演化，并希望通过这种演化逐渐消除局部的复杂性。Ricci flow 的思想与热方程具有深刻的类比。曲率在演化过程中发生扩散，而几何结构则逐渐趋向更加规则的状态。然而，与普通热方程不同的是，Ricci flow 会产生奇异性。

因此，真正困难的问题并不是证明 Ricci flow 存在，而是理解在 $t \nearrow T$ 时流形究竟发生了什么，以及在奇异时刻之后如何继续理解这个几何对象。Perelman 对 Ricci flow 的工作最终完成了这一计划。通过 entropy、non-collapsing、surgery 以及奇异性分析等一系列工具，三维流形可以被分解成具有高度受控几何结构的部分，从而最终解决了 Poincar\'e 猜想以及更一般的 Thurston 几何化猜想。

然而，从今天的角度来看，Ricci flow 带来的最重要遗产或许并不是 Poincar\'e 猜想本身。真正重要的是：

\begin{quote}
我们开始系统地研究一个光滑几何对象在退化之后究竟会变成什么。
\end{quote}

当一列黎曼流形 $(M_i,g_i)$ 在适当的曲率和几何条件下收敛时，即使每一个 $M_i$ 都是光滑流形，它们的极限 $(X,d)$ 也未必仍然是光滑流形。奇异空间并不是人为构造出来的反例，而是自然几何演化和几何极限中不可避免的对象。于是现代几何分析发生了一次根本性的转变：研究光滑空间 $\longrightarrow$ 研究光滑空间的极限。

\section{从 Ricci flow 到 Ricci limit spaces}

在 Ricci curvature 得到统一下界的条件下，Gromov 的紧性理论给出了度量空间层面的预紧性，而 Cheeger--Colding 则进一步系统研究了这些极限空间的内部结构。考虑 $\operatorname{Ric}_{g_i} \geq K g_i$。即使每个 $(M_i,g_i)$ 都是 $n$ 维光滑黎曼流形，其 Gromov--Hausdorff 极限 $(X,d)$ 仍然可能出现奇异点。然而，这些奇异空间并不是完全任意的。

Cheeger--Colding 理论表明，Ricci 曲率下界会以一种极其强的方式限制极限空间的结构。Regular set、singular set、tangent cones、splitting maps、almost splitting theorem 以及 volume convergence 等理论逐渐形成了一套完整的 Ricci limit space 理论。这里发生了一件非常重要的事情：

\begin{quote}
奇异性本身开始成为几何对象。
\end{quote}

这意味着低正则性并不等于完全失去结构。恰恰相反，曲率条件本身能够在极限空间中产生新的正则性和刚性。然而，Cheeger--Colding 理论仍然从黎曼流形序列出发。于是一个更加大胆的问题自然出现：

\begin{quote}
如果我们完全忘记极限空间来自黎曼流形这一事实，能否仅仅从度量和测度本身定义 Ricci curvature？
\end{quote}

\section{从 Ricci curvature 到 optimal transport}

Lott--Sturm--Villani 给出了一个革命性的答案。他们发现，Ricci curvature 的下界可以通过概率测度空间上的 Wasserstein 几何来刻画。对于一个测度度量空间 $(X,d,m)$，考虑概率测度空间 $\mathcal{P}_2(X)$ 以及 Wasserstein distance $W_2$。Ricci curvature 的下界可以被转化为熵泛函沿 Wasserstein geodesic 的凸性。于是：Ricci curvature $\longrightarrow$ optimal transport。

Ricci curvature 从一个依赖二阶微分结构的局部对象，变成了一个可以通过距离、测度和概率分布之间的运输结构描述的对象。由此产生了 Curvature-Dimension condition，即 $CD(K,N)$。这一转变的重要性远远超过了给 Ricci curvature 找到一个新的定义。它意味着：

\begin{quote}
一个原本属于光滑微分几何的结构，可能存在一个完全不依赖微分结构的本质描述。
\end{quote}

这为真正意义上的奇异 Ricci 几何打开了大门。

\section{RCD：从曲率到能量}

然而，$CD(K,N)$ 空间仍然过于宽泛。特别是，$CD$ 条件并不能排除 Finsler geometry。因此，如果我们希望得到真正意义上黎曼几何的奇异推广，还需要加入能够刻画 Hilbertian structure 的条件。这便产生了 RCD 空间。粗略而言，
\[
    RCD(K,N) = CD(K,N) + \text{infinitesimal Hilbertianity}.
\]
其中 infinitesimal Hilbertianity 可以通过 Cheeger energy
\[
    \operatorname{Ch}(f) = \frac{1}{2} \int_X |Df|^2 \, dm
\]
的二次性来表达。这一步尤其重要。

我们并没有直接把黎曼度量重新放回定义中。相反，我们要求能量结构具有 Hilbert 空间的二次结构，从而让梯度、Laplace operator、heat flow 等分析对象重新出现。于是出现了一条非常漂亮的路线：
\[
    \text{Riemannian manifolds} \longrightarrow \text{Ricci limit spaces} \longrightarrow \text{RCD spaces}.
\]
RCD 理论因此成为一种非常成功的 synthetic Ricci geometry。它不仅包含大量 Ricci limit spaces，而且在适当的测度 Gromov--Hausdorff 收敛下具有良好的稳定性。Heat flow、Cheeger energy、Laplacian、Bochner inequality、Sobolev calculus、rectifiability、splitting theorem 以及大量经典几何分析结果，都可以在这一框架中得到相应版本。到这里，一切看起来都非常成功。然而，正是在这里，一个更加危险的问题开始出现。

\section{迷失的二十年}

过去二十年的 RCD 理论取得了令人惊叹的技术成果。问题恰恰不是 RCD 做得不够多。问题是：

\begin{quote}
RCD 做得太成功了。
\end{quote}

它建立了一套非常强大的工具，使得我们能够在没有光滑结构的空间上重新定义梯度、Hessian、Laplacian、heat flow 以及大量经典分析对象。于是研究逐渐形成了一种非常稳定的模式：经典黎曼几何 $\longrightarrow$ 去掉光滑性 $\longrightarrow$ 寻找 synthetic formulation $\longrightarrow$ 证明经典定理仍然成立 $\longrightarrow$ 建立更多 calculus $\longrightarrow$ 证明更多经典定理。这个循环产生了大量漂亮而深刻的数学。但是，它也可能正是过去二十年低正则性几何最值得警惕的问题。

\subsection{我们究竟在研究奇异空间，还是在恢复黎曼流形？}

RCD 空间可以具有真正的奇异性。它可能不是流形，tangent cone 可以出现复杂结构，局部维数可以发生变化，甚至微分结构本身也只能以弱意义存在。然而，大量研究的基本目标却是证明：RCD space $\approx$ Riemannian geometry almost everywhere。Rectifiability theory 是最典型的例子。我们不断寻找 coordinate charts、Euclidean tangent cones、splitting maps 以及弱微分结构，最终证明空间在几乎处处意义下具有良好的欧氏性质。这当然是重要的。但是这里存在一个更根本的问题：

\begin{quote}
如果我们的最终目标始终是证明一个奇异空间“几乎处处仍然是黎曼的”，那么我们是否真的在研究奇异性？
\end{quote}

我们真正应该问的问题也许恰恰相反：“singularity itself carries what geometry?” 而不仅仅是 “how small is the singular set?”

\subsection{从恢复经典定理到创造新定理}

过去二十年的另一个显著特征，是大量经典黎曼几何定理被推广到 RCD 空间。Splitting theorem、volume cone implies metric cone、maximum diameter theorem、Obata theorem、isoperimetric inequalities、Sobolev inequalities、spectral convergence 等，都在 RCD 框架中获得了相应的版本。这些工作当然具有巨大的价值。但我们需要区分两个完全不同的问题：哪些经典结构能够在极限中保留下来？和：极限本身创造了哪些新的结构？

前一个问题已经得到了极其丰富的研究。后一个问题却远没有得到同样程度的关注。这导致了一种奇怪的局面：

\begin{quote}
我们越来越了解 RCD 空间为什么仍然像黎曼流形，却没有同样清楚地理解它们为什么不再是黎曼流形。
\end{quote}

换句话说，低正则性理论中的“奇异”经常被当作一个需要控制的误差，而不是一个需要解释的数学现象。

\subsection{稳定性是不是成为了终点？}

RCD 理论最令人惊叹的性质之一当然是稳定性。若 $(X_i,d_i,m_i) \longrightarrow (X,d,m)$ 在适当意义下收敛，并且每个 $X_i$ 都满足统一的 RCD 条件，那么极限空间仍然保留相应的 RCD 结构。这建立了一个非常强的闭性原理：Ricci lower bound + geometric convergence $\Longrightarrow$ Ricci lower bound on the limit。然而稳定性回答的是：经典结构退化以后还剩下什么？它并没有回答：退化以后产生了什么？

这两个问题不能混为一谈。稳定性当然是一个理论必须具有的重要性质。但如果一个理论只关心结构是否在极限中保持，那么它很容易把极限中的新现象重新解释成原有结构的弱极限。于是出现了一个值得警惕的可能：stability may preserve structure while hiding novelty。

\subsection{工具开始决定问题}

RCD 理论逐渐形成了一套非常成熟的技术闭环：optimal transport $\longrightarrow$ entropy convexity $\longrightarrow$ Cheeger energy $\longrightarrow$ heat flow $\longrightarrow$ Bochner inequality $\longrightarrow$ second-order calculus。这套理论极其漂亮。但它同时产生了一种方法论上的吸引力。一旦一个问题能够被写成：energy + convexity + gradient flow，那么 RCD machinery 就可以自然地介入。于是研究者很容易继续寻找那些能够被这套 machinery 处理的问题。这可能导致一种反常现象：

\begin{quote}
我们拥有越来越强大的工具，但研究的问题本身却越来越接近这些工具能够处理的问题。
\end{quote}

最终，methodology begins to determine the geometry。这或许才是所谓“迷失的二十年”真正值得警惕的地方。

\section{一个真正困难的问题}

因此，我们应该重新提出一个更加尖锐的问题：

\begin{quote}
RCD 空间作为黎曼流形、加权黎曼流形以及 Ricci limit spaces 的统一扩张，究竟有没有产生一种不能从经典黎曼几何直接看见的新几何结构？
\end{quote}

这里并不是说 RCD 空间没有新现象。恰恰相反，它们拥有大量经典流形中不存在的奇异行为。问题是这些现象通常被组织成：如何恢复 regularity、如何控制 singularity、如何证明 stability，以及如何推广 classical theorem。而真正缺少的是另一种研究方向：从 singularity 出发，而不是消除 singularity。也就是说，我们需要从 What survives the limit? 转向 What is created by the limit? 这也许是过去二十年 RCD 理论最重要而又最少被正面提出的问题。

\section{为什么 Lorentzian geometry 是一个自然的实验场}

如果我们希望寻找低正则性理论真正可能产生的新结构，Lorentzian geometry 是一个极其自然的实验场。黎曼几何与 Lorentzian geometry 都以 metric 为核心，但二者的几何本质并不相同。在 Lorentzian geometry 中，我们不仅有 metric，还有 causal structure 和 time separation。记 time separation 为 $\tau(x,y)$。它不是普通的距离函数。它受到 causal relation $x \leq y$ 的控制。因此 Lorentzian geometry 至少同时包含：metric + causal structure + time separation。

这意味着把 RCD 的思想直接搬到 Lorentzian 世界，从一开始就存在一个根本性的风险：

\begin{quote}
我们可能成功推广了 Ricci curvature，却没有真正推广 Lorentzian geometry。
\end{quote}

如果我们只研究 Lorentzian analogue of $CD(K,N)$，那么最容易得到的是体积变化率、Jacobian estimates 以及 optimal transport 中的熵凸性。这些结果当然重要。但它们主要描述的是 volume distortion。而 Lorentzian geometry 最独特的结构却是：causality and time separation。如果一个理论不能真正解释 $\tau(x,y)$、$J^+(x)$、$J^-(x)$ 以及 causal geodesics，那么我们究竟是在构造 Lorentzian geometry 的低正则性版本，还是仅仅构造了一个带有 Lorentzian 味道的 CD 空间？这是两个完全不同的问题。

\section{更加危险的问题：PDE}

问题甚至比 causal structure 更深。RCD 理论之所以能够取得巨大成功，一个关键原因是黎曼几何中的 PDE 与能量结构具有极其紧密的联系。在光滑黎曼流形上，$\Delta f = \operatorname{div}(\nabla f)$ 是椭圆算子。相应的 Bochner identity 为
\[
\frac{1}{2}\Delta |\nabla f|^2 = |\operatorname{Hess}f|^2 + \langle \nabla f, \nabla\Delta f \rangle + \operatorname{Ric}(\nabla f,\nabla f).
\]
它把 PDE + energy + Ricci curvature 联系在了一起。然而 Lorentzian geometry 中对应的 d'Alembert operator $\Box_g$ 却是双曲型算子。因此 $\Delta_g$ 与 $\Box_g$ 在 PDE 理论中的性质存在根本差异。

Riemannian geometry 中的基本分析结构是：elliptic PDE, heat flow, energy minimization。而 Lorentzian geometry 中则是：hyperbolic PDE, wave propagation, causal evolution。因此：heat flow $\neq$ wave propagation，gradient flow $\neq$ causal evolution。如果我们仅仅把 RCD 中的 Cheeger energy $\longrightarrow$ heat flow $\longrightarrow$ Laplacian 这一套结构搬到 Lorentzian geometry，那么我们很可能从根本上错过 Lorentzian geometry 最重要的部分。Lorentzian geometry 中的信息不是简单地向空间中扩散。它具有方向。信息沿 causal cones 传播。因此，未来的 synthetic Lorentzian geometry 很可能不能仅仅建立在 metric measure structure 和 convexity 之上。它必须能够同时编码：causality + time separation + hyperbolic PDE + Ricci curvature。

\section{Einstein equation：真正的终点}

Lorentzian geometry 还有一个 Riemannian geometry 无法回避的问题：Einstein field equation。在光滑 Lorentzian manifold 上，
\[
    \operatorname{Ric}_{\mu\nu} - \frac{1}{2}R g_{\mu\nu} + \Lambda g_{\mu\nu} = 8\pi T_{\mu\nu}.
\]
这里 Ricci curvature 并不是一个孤立的曲率条件。它是一个动力学方程的一部分。因此，如果我们希望建立真正的 synthetic spacetime theory，仅仅定义一个 Ricci curvature lower bound 是远远不够的。我们真正需要问：

\begin{quote}
在没有经典微分结构的情况下，Einstein equation 是否仍然具有一个自然的意义？
\end{quote}

这与 RCD 理论中的问题有本质区别. RCD 主要回答：Ricci lower bound 如何在奇异空间中表达。而 Lorentzian geometry 最终需要回答：Einstein dynamics 如何在奇异空间中表达。这两者之间存在巨大的鸿沟。

\section{从“推广定理”到“发现结构”}

因此，未来二十年的问题不应该只是：如何把更多 Lorentzian geometry 的经典定理推广到 LCD spaces？而应该是：低正则性究竟迫使我们发现什么新的几何结构？第一种研究方式是：classical geometry $\longrightarrow$ weak formulation $\longrightarrow$ synthetic theorem。第二种研究方式则是：classical geometry $\longrightarrow$ loss of regularity $\longrightarrow$ new phenomenon $\longrightarrow$ new structure。

两者的差异并不是技术路线上的差异。它们实际上代表了两种完全不同的数学哲学。过去二十年，我们已经非常成功地回答了：光滑性消失以后，经典黎曼几何还能保留下来多少？但未来真正值得回答的问题也许应该是：光滑性消失以后，什么东西第一次出现了？

\section{未来二十年的问题}

从这个角度来看，未来低正则性 Lorentzian geometry 至少存在几个基本问题。

\subsection{Ricci curvature 应该如何被重新理解？}

如果底层空间没有微分结构，那么 $\operatorname{Ric}$ 不再是一个可以逐点定义的 tensor。仅仅通过 entropy convexity 定义 Ricci lower bound，是否已经包含 Einstein equation 所需要的全部信息？还是应该存在一种同时看到 $\operatorname{Ric}$ and causal geometry 的新对象？

\subsection{Time separation 是否应该成为基本变量？}

在 Lorentzian geometry 中，$\tau(x,y)$ 可能比一个抽象的 Lorentzian metric 更加接近真正的基本变量。因此，一个真正的 synthetic spacetime 也许应该从 $(X,d,m)$ 走向某种 $(X,\tau,m,\leq)$。这意味着理论将从 metric measure geometry 走向 causal measure geometry。

\subsection{是否存在 Lorentzian analogue of Cheeger energy？}

RCD 的 infinitesimal Hilbertianity 最终由 Cheeger energy 的二次性体现。那么 Lorentzian 世界中对应的能量结构是什么？一个自然的候选并不一定是正定的 Dirichlet form，因为 Lorentzian metric 本身具有不定号。因此，我们也许需要某种完全不同的能量结构：indefinite energy，或者更加接近 Hamiltonian mechanics 的结构。

\subsection{Einstein equation 能否成为 synthetic equation？}

最终最重要的问题或许是：Can Einstein's equation survive without smoothness? 如果答案仅仅是：可以定义一个 synthetic Ricci curvature lower bound。那么这仍然只是 Ricci geometry 的推广。真正令人兴奋的答案应该是：Einstein spacetime $\longrightarrow$ synthetic Einstein spacetime。如果能够做到这一点，那么我们得到的就不再是 Lorentzian geometry 的一个弱版本，而可能是一种真正新的几何理论。

\section{结语：从二十年的继承到二十年的创造}

从 Poincar\'e 猜想到 Ricci flow，从 Ricci flow 到 Ricci limit spaces，从 Cheeger--Colding 到 Lott--Sturm--Villani，再到 RCD spaces，现代几何学花费了一个世纪，逐渐学会了如何在没有光滑结构的情况下继续讨论 Ricci curvature。这条道路告诉我们：smoothness is not the essence of geometry。但它也留下了一个更加困难的问题：if smoothness is not essential, what is?

过去二十年的 RCD 理论，已经非常成功地回答了：经典黎曼几何能够保留下来多少？但我们仍然没有真正回答：低正则性本身创造了什么？如果过去的主要任务是不断证明 RCD space $\approx$ Riemannian geometry almost everywhere，那么下一个二十年真正值得研究的问题也许恰恰相反：RCD space $\not\approx$ Riemannian geometry。

我们需要寻找的不是：哪里仍然是欧氏空间？而是：哪里开始不再是欧氏空间？不是：哪些经典定理仍然成立？而是：哪些新的定理只有在奇异空间中才可能出现？不是：如何恢复微分结构？而是：没有微分结构以后，什么结构取代了它？

这也是为什么 Lorentzian geometry 如此重要。因为在 Lorentzian 世界中，我们无法简单地把 Riemannian theory 的每一个工具重新建立一次。因果性不会被 rectifiability 消除。Time separation 不会被 Cheeger energy 替代。Hyperbolic propagation 不会因为拥有一个良好的 heat flow 而消失。Einstein equation 也不会因为存在一个 synthetic Ricci lower bound 而自动出现。

因此，Lorentzian geometry 可能迫使低正则性理论真正完成它尚未完成的任务：从恢复经典结构，走向发现新结构。如果说过去二十年是在回答：光滑性消失以后，黎曼几何还能保留多少？那么接下来的二十年，也许应该开始回答：光滑性消失以后，几何究竟会变成什么？

这或许才是低正则性几何真正尚未开始的部分。

\end{document}
